% Modernised reading — fonds Alexandre Grothendieck, shelfmark 120, pages 1-24.
%
% This is an INTERPRETATION, not a transcription. Everything the critical
% apparatus carried moves into footnotes. In any dispute of substance the
% transcription governs: whoever cites Grothendieck must cite
% batch-01.fr.tex and batch-02.fr.tex.
%
% Pass: Opus 5.5 (claude-opus-5-5), 2026-10-03 — first pass, unchecked against the pages by a human.
% The folder was read whole: two batches, pages 1-24, both transcribed. Page 1
% is the cover; it carries no \page in the transcription.
%
% Notation fixed for the whole document. R a set, R_0 its tame points;
% \mathcal{M}_n (n >= 1) and \mathcal{M}_I (I finite) the tame parts;
% \Sigma = Esp_{\mathcal{M}_*} the category of tame spaces, |X| the underlying
% set; \Sigma^{ind} for the category of strict ind-tame spaces, which the page
% writes « Esp_M » with an overwritten index (p. 14); gothic letters for
% ind-tame spaces, as on pp. 15-16. The page uses the label (M5) twice: here
% (M5) is R_0 != Ø (p. 6) and (M5') is stability under finite unions (p. 9);
% every later « M5 » of the pages (pp. 12, 14) is (M5'). The axioms are
% otherwise numbered as he numbers them, (M0)-(M11).
%
% Substantive departures from the pages, each footnoted where it occurs:
% — p. 2: « les \mathfrak{S}_n stables par intersection » read \mathcal{M}_n;
%   « pour avoir M3 il suffit de le prouver pour u surjectif » read u
%   injective (u^* surjective), as the page's own c_1) of p. 3 has it.
% — p. 3: the reduction a)-c) replaces M1-M3 only together with M4.
% — p. 4: « {x} \in R^n » read {x} \in \mathcal{M}_n.
% — p. 5: the margin « via M2 » — our proof of a)-c) uses M1, M3, M4.
% — p. 6: the final object exists without (M5) once I = Ø is admitted (p. 4,
%   p. 5 margin); with I >= 1 only, (M5) is what supplies it.
% — p. 7: the characterisation by (\Sigma, R, |.|) — condition b) illegible;
%   we state only the recovery of \mathcal{M}_* from the triple.
% — p. 8: the discrete-structure corollary needs (M5') (counterexample:
%   example a) of p. 5); the colimit Proposition of p. 8 also uses (M5') in
%   its proof (p. 11), except for coequalisers.
% — p. 10: the functions f_alpha, g_alpha of the criterion must be tame
%   (supplied); « X_1 -> R^I » read R^A.
% — p. 14: « (M8) impliqué par M5 » is false for either M5 (counterexample);
%   M8 follows from M0.
% — p. 14: finite sums of ind-tame spaces under « R_0 != Ø »: our
%   construction needs two distinct tame points.
% — p. 15: « dans Esp_{M_*} » read in \Sigma^{ind}.
% — p. 17: the colimit topology induces on each tame part its own topology
%   already without countability (argument in a footnote); « dénombrable »
%   is an uncertain reading.
% — p. 20: the finite limits of \Sigma need M3 (or M2) besides M1 and M4.
% — p. 22: e_alpha \in R corrected to e_alpha \in R_0, as M6 has it.
% — p. 23: « sépare les points de X » read X \smallsetminus Y.
%
% What the folder announces and never establishes: the converse of the
% characterisation of p. 7; the existence of quotients and amalgamated sums
% of ind-tame spaces (p. 15); the stability of locally tame spaces under
% finite limits (p. 19); the « Dém : » of p. 23 (left empty; the argument
% stands above it); the whole programme « À prouver » of p. 24 (Hardt-type
% triviality). The \mathbb{D} opening p. 21 is not explained by the folder.

\documentclass[11pt,a4paper]{article}
\input{../preamble/grothendieck}

\folder{120}
\pages{1}{24}
\dating{[à partir de 1973-à partir de 1978]}
\watermark{Édition de démonstration\\interprétation personnelle de l'œuvre}
\foldertitle{Topologie modérée : notes manuscrites (s.d.) — lecture modernisée du dossier entier}

\begin{document}

\section*{Résumé}

\begin{resume}
Ces vingt-quatre pages cherchent à dire, par un petit nombre de règles,
quelles figures de l'espace méritent d'être appelées « sages ».

La topologie générale admet des objets que personne ne dessine : une courbe
qui remplit un carré, un ensemble dense et de complémentaire dense, un compact
qui a la forme d'une poussière. Ce ne sont pas des erreurs, mais ce n'est pas
de la géométrie. La \emph{topologie modérée} que Grothendieck appelle de ses
vœux fait le choix inverse : décider d'avance quelles parties sont admises —
les figures définies par des polynômes, par exemple, et ce qu'on en tire par
des opérations raisonnables — et ne jamais sortir de cette famille. On perd les
monstres ; on espère gagner des théorèmes de finitude.

Le dossier ne commence pas par les figures, mais par les \emph{règles}. On se
donne un ensemble $R$ — la droite, pensée abstraitement — et, pour chaque
$n$, une famille de parties de $R^n$ déclarées modérées. Les règles disent
seulement ce qu'on a le droit de faire : permuter ou dupliquer des
coordonnées, projeter, couper par une condition, former des produits. Une
image aide : c'est un jeu de construction. On fournit quelques pièces et
quelques façons de les assembler, et l'on ne demande rien d'autre ; tout ce
qu'on peut monter est admis, rien d'autre ne l'est. Les deux premières
sections montrent que ces règles suffisent à fabriquer une catégorie
d'« espaces modérés » où l'on peut former produits, équations et
sous-espaces exactement comme dans les ensembles.

La question qui occupe ensuite presque tout le dossier est celle du
\emph{recollement} : peut-on coller deux espaces modérés le long d'une partie
commune, ou écraser une partie en un point, sans sortir du jeu ? Projeter et
couper ne posent pas de problème ; recoller en pose un, parce qu'il faut
trouver, dans un $R^n$, une place où loger le résultat. Grothendieck donne un
critère, puis deux axiomes qui le rendent applicable : toute partie modérée
est le lieu où un nombre fini de fonctions modérées prennent des valeurs
fixées, et toute fonction modérée sur une partie se prolonge à l'espace
entier. Moyennant quoi les recollements existent, et la preuve est une
manière de ranger les deux morceaux dans des étages différents d'un même
$R^n$ pour qu'ils ne se rencontrent que là où ils doivent.

Viennent ensuite trois greffes sur ce tronc. Les \emph{ind-espaces}
autorisent les objets non bornés, comme la droite entière, réunions
croissantes de pièces modérées. La \emph{topologie} demande que les pièces
soient compactes ; les applications modérées deviennent alors automatiquement
continues, et l'on peut parler de ce qui est modéré au voisinage d'un point.
La \emph{structure de corps} demande que l'addition et la multiplication
soient modérées ; les fonctions modérées forment alors un anneau, les points
modérés un corps — pour les polynômes à coefficients rationnels, ce serait le
corps des nombres algébriques réels —, et une fonction modérée qui ne
s'annule pas a un inverse modéré.

La dernière page n'est plus une construction mais un programme, marqué
« À prouver » : au-dessus d'une base modérée, toute application modérée
devrait être un produit trivial hors d'une partie mince, morceau par morceau.

Ces pages anticipent ce qui s'est appelé, quelques années plus tard, les
\emph{structures o-minimales} de la théorie des modèles, dont le livre de
référence porte le titre que Grothendieck avait choisi ; le programme de la
dernière page est, dans ce cadre, le théorème de trivialité de Hardt. La
feuille de garde porte, au crayon, « Autour de l'Éloge » : c'est le même
programme que l'\emph{Esquisse d'un programme} défendra en 1984.

\keywords{tame topology, o-minimal structure, definable sets, semi-algebraic sets, concrete category, transport of structure, finite limits, pushout along monomorphisms, quotient by a closed subset, definable Tietze extension, strict ind-object, colimit topology, sheaf of definable functions, field of definition, Hardt triviality}
\end{resume}

\pagerange{1}{24}
\section*{Le dossier, sa charpente et ses notations}

Vingt-quatre pages, toutes de sa main, d'une même encre et d'une rédaction
rapide. La page 1 est une chemise rose ; les pages 2 à 23 sont un exposé
axiomatique, divisé par lui en six sections numérotées en chiffres romains
cerclés ; la page 24 est un programme. Il numérote lui-même un feuillet sur
deux environ, de 1 (page 2) à 11 (page 22), sans lacune : le dossier est une
seule rédaction, non des fragments réunis.

Les stations de l'argument :

\begin{itemize}
\item[(I)] les \emph{modèles $R$-modérés} : les axiomes (M1) à (M4) sur des
  familles $\mathcal{M}_n$ de parties de $R^n$, leurs réductions, les cas
  triviaux, les points modérés (pages 2 à 5) ;
\item[(II)] les \emph{espaces modérés} : la catégorie $\Sigma$, ses produits et
  limites finies, la droite $R$ comme objet de $\Sigma$ (pages 5 à 7) ;
\item[(III)] les \emph{limites inductives finies} : critère d'existence des
  sommes amalgamées, quotients, les axiomes (M6) et (M7), et l'existence des
  sommes amalgamées le long de plongements (pages 7 à 13) ;
\item[(IV)] les \emph{ind-espaces modérés stricts} (pages 14 à 16) ;
\item[(V)] la \emph{topologie} : compacité des parties modérées, voisinages
  modérés, espaces localement modérés (pages 16 à 19) ;
\item[(VI)] la \emph{structure de corps} : l'anneau des fonctions modérées, le
  corps $R_0$, et le retour sur (M6) qu'elle permet (pages 20 à 23) ;
\end{itemize}
puis, page 24, le programme « À prouver ».

Le tronc est (I)--(III) : construire, à partir de règles portant seulement sur
des parties de $R^n$, une catégorie où les limites finies et certaines
limites inductives finies existent et se calculent comme dans les ensembles.
Les sections (IV) à (VI) sont des greffes indépendantes, chacune avec ses
axiomes propres ; (VI) commence d'ailleurs par « Oublions les structures
topologiques ». La section (VI) revient cependant sur (III) : elle montre que
la seconde moitié de l'axiome (M6) est conséquence de la première dès que $R$
est un corps.

\textbf{Notations, valables pour tout le dossier.} $R$ est un ensemble
quelconque, sans structure, jusqu'à la section (V). Pour $n \geqslant 1$,
$\mathcal{M}_n$ est un ensemble de parties de $R^n$ ; ses éléments sont les
\emph{parties modérées}. Pour une application $u : [1,n] \to [1,m]$, on note
$u^* : R^m \to R^n$, $x \mapsto x \circ u$, l'application qu'elle induit sur
les coordonnées : une projection quand $u$ est injective, un plongement
diagonal quand $u$ est surjective. $R_0 \subset R$ est l'ensemble des
\emph{points modérés}, ceux dont le singleton est modéré. $\Sigma =
\mathrm{Esp}_{\mathcal{M}_*}$ est la catégorie des espaces modérés et $X
\mapsto |X|$ son foncteur d'oubli ; $\Sigma^{\mathrm{ind}}$ est celle des
ind-espaces modérés stricts, dont les objets sont notés en lettres gothiques
$\mathfrak{X}, \mathfrak{Y}, \mathfrak{Z}$, comme aux pages 15 et 16.

Les axiomes gardent ses numéros, à une exception près. Il emploie deux fois
l'étiquette (M5) : page 6 pour $R_0 \neq \emptyset$, page 9 pour la stabilité
par réunions finies. On garde (M5) pour la première et l'on écrit (M5$'$) pour
la seconde ; toutes les mentions ultérieures de « M5 » (pages 12 et 14)
renvoient à la seconde. Récapitulons, pour s'y retrouver :

\begin{itemize}
\item[(M0)] $R \in \mathcal{M}_1$ (facultatif, page 6) ;
\item[(M1)--(M4)] images directes par les $u^*$, images réciproques par les
  plongements diagonaux, intersections avec des images réciproques, produits
  (page 2) ;
\item[(M5)] $R_0 \neq \emptyset$ (page 6) ; \quad (M5$'$) stabilité par
  réunions finies (page 9) ;
\item[(M6)] toute partie modérée d'une partie modérée est une fibre d'un
  nombre fini de fonctions modérées ; (M7) prolongement des fonctions modérées
  (page 12) ;
\item[(M8)] $R$ est réunion de ses parties modérées (page 14) ;
\item[(M9)] $R$ séparé et parties modérées compactes ; (M10) voisinages
  modérés (pages 16 et 17) ;
\item[(M11)] axiomes de corps (page 20).
\end{itemize}

Un mot sur le mot « modéré ». Une \emph{partie} modérée est un élément d'un
$\mathcal{M}_n$ ; une \emph{application} modérée est une application dont le
graphe est une partie modérée ; une \emph{fonction} modérée est une application
modérée à valeurs dans $R$. Ce sont trois emplois d'une même idée — le graphe
ramène tout aux parties — et c'est le premier geste du dossier (page 5).

Le dossier voisin 121, daté de 1974, reprend aux pages 53 à 72 une démarche
très proche, mais en prenant pour pièces de base les parties des cubes
$[0,1]^n$ plutôt que celles des $R^n$.

\pagerange{1}{1}
\section*{La chemise (page 1)}

La chemise rose porte à l'encre, de sa main, « Topologie modérée », le premier
mot écrit par-dessus un mot biffé, peut-être « Structures », et la finale de
« modérées » corrigée en « modérée ».\footnote{Lecture de la transcription,
qui donne « Structures » comme douteux. Le titre initial aurait donc été
« Structures modérées », ce qui décrit mieux les pages 2 à 23 que le titre
retenu : il s'agit de structures sur des ensembles, la topologie n'arrivant
qu'à la page 16.} Au-dessous, au crayon très pâle et peut-être d'une autre
main, « 1978 » suivi d'un mot illisible, et « (Autour de l'Éloge) ».\footnote{La
lecture « Autour de l'Éloge » est douteuse. Si elle est juste, elle rattache
ces pages à un « éloge » de la topologie modérée, ce que l'\emph{Esquisse d'un
programme} (1984) développe en effet ; l'édition n'en tire pas de datation.
La cote de l'inventaire, « [à partir de 1973-à partir de 1978] », est
reproduite telle quelle dans le titre.}

\pagerange{2}{5}
\section*{(I) Modèles $R$-modérés (pages 2 à 5)}

\subsection*{Les axiomes}

On se donne, pour tout $n \geqslant 1$, un ensemble $\mathcal{M}_n$ de parties
de $R^n$, et l'on demande :

\begin{itemize}
\item[(M1)] pour toute application $u : [1,n] \to [1,m]$ et tout $A \in
  \mathcal{M}_m$, on a $u^*(A) \in \mathcal{M}_n$ ;
\item[(M2)] pour toute $u : [1,n] \to [1,m]$ \emph{surjective} — de sorte que
  $u^* : R^m \hookrightarrow R^n$ est injective — et tout $B \in
  \mathcal{M}_n$, on a $(u^*)^{-1}(B) \in \mathcal{M}_m$ ;
\item[(M3)] pour toute $u : [1,n] \to [1,m]$, $A \in \mathcal{M}_m$ et $B \in
  \mathcal{M}_n$, on a $(u^*)^{-1}(B) \cap A \in \mathcal{M}_m$ ;
\item[(M4)] pour $A \in \mathcal{M}_m$ et $B \in \mathcal{M}_n$, on a $A
  \times B \in \mathcal{M}_{m+n}$, via $R^m \times R^n \simeq R^{m+n}$.
\end{itemize}

En clair : (M1) dit que les parties modérées sont stables par permutation et
duplication des coordonnées et \emph{par projection} ; (M2), qu'on peut
couper une partie modérée par une diagonale ; (M3), qu'on peut la couper par
la condition « telles coordonnées tombent dans telle partie modérée » ; (M4),
qu'on peut faire des produits. La stabilité par projection est la règle
décisive : c'est, dans le cas des parties semi-algébriques, le théorème de
Tarski–Seidenberg, et c'est elle qui fait qu'une partie définie par « il
existe $y$ tel que \ldots » reste modérée.\footnote{Le rapprochement avec
Tarski–Seidenberg est de l'édition ; la page ne nomme aucun exemple
géométrique. Noter aussi ce qui \emph{manque} à la liste : ni réunions ni
complémentaires. Les parties modérées ne forment pas une algèbre de Boole,
contrairement aux ensembles définissables d'une structure o-minimale ; la
réunion n'arrive qu'avec (M5$'$) page 9, et le complémentaire jamais, ce qui
s'accorde avec la compacité demandée en (M9) page 16 : il pense à des
parties fermées bornées.}

Par (M3) appliqué à $u = \mathrm{id}$, les $\mathcal{M}_n$ sont stables par
intersection de deux parties.

\subsection*{Réductions}

Comme (M1) est stable par composition et que toute application
$[1,n] \to [1,m]$ est composée d'une surjection et d'une injection, (M1)
revient à le demander pour les permutations, pour les « opérations de face »
($u : [1,n] \hookrightarrow [1,n+1]$ injective, $u^*$ une projection
$R^{n+1} \to R^n$) et pour les « opérations de dégénérescence » ($u :
[1,n+1] \to [1,n]$ surjective, $u^*$ un plongement diagonal $R^n
\hookrightarrow R^{n+1}$) — et même, à permutation près, pour la seule
projection $(x_1, \dots, x_{n+1}) \mapsto (x_1, \dots, x_n)$ et le seul
plongement $(x_1, \dots, x_n) \mapsto (x_1, \dots, x_n, x_n)$. De même (M2) se
réduit à ce plongement diagonal standard.

(M1), (M2) et (M4) entraînent la stabilité des $\mathcal{M}_n$ par
intersection, sans (M3) : $A \cap B$ est l'image réciproque de $A \times B$
par la diagonale $R^n \hookrightarrow R^{2n}$.\footnote{La page écrit « les
$\mathfrak{S}_n$ sont stables par intersection » et « $A, B \in
\mathfrak{S}_n \Rightarrow A \cap B \in \mathfrak{S}_n$ », avec le sigle du
groupe symétrique employé trois lignes plus haut ; c'est $\mathcal{M}_n$ qui
est entendu.} Il en résulte que (M3) n'a à être demandé que lorsque $u^*$ est
une projection, et même pour la seule projection standard $R^{n+1} \to
R^n$.\footnote{La page dit « pour $u$ surjectif ». C'est le cas inverse qui
reste à traiter : si $u$ est surjective, $(u^*)^{-1}(B)$ est déjà modérée par
(M2) et son intersection avec $A$ l'est par stabilité par intersection ; le
cas non trivial est celui où $u$ est \emph{injective} et $u^*$ une projection.
C'est bien ce cas, et lui seul, que retient le c$_1$) du résumé de la page 3,
avec $\beta_n : [1,n] \hookrightarrow [1,n+1]$. Le passage d'une projection
quelconque à la projection standard se fait par récurrence, en écrivant, pour
$p = p_1 \circ p_2$, $A \cap p^{-1}(B) = A \cap p_2^{-1}\bigl(p_2(A) \cap
p_1^{-1}(B)\bigr)$, où $p_2(A)$ est modérée par (M1).}

En résumé, en présence de (M4), les axiomes (M1) à (M3) équivalent à :

\begin{itemize}
\item[a)] stabilité de chaque $\mathcal{M}_n$ par les permutations des
  coordonnées ;
\item[b)] pour le plongement diagonal $\alpha_n^* : R^n \hookrightarrow
  R^{n+1}$, $(x_1, \dots, x_n) \mapsto (x_1, \dots, x_n, x_n)$ : stabilité par
  image directe (b$_1$) et par image réciproque (b$_2$) ;
\item[c)] pour la projection $\beta_n^* : R^{n+1} \to R^n$ : stabilité de $A
  \cap (\beta_n^*)^{-1}(B)$ (c$_1$) et stabilité par image directe (c$_2$).\footnote{La page dit « M1, M2, M3 peuvent être remplacés par » a), b), c),
sans mentionner (M4). Il faut (M4), ou (M3) en entier, pour obtenir la
stabilité par intersection dont la réduction de (M3) se sert ; d'où la
précision « en présence de (M4) ». En b$_1$), l'exposant de $\alpha_n$ est mal
formé sur la page ; c'est bien une image directe.}
\end{itemize}

Enfin (M1) permet de définir $\mathcal{M}_I \subset \mathfrak{P}(R^I)$ pour
tout ensemble fini non vide $I$, par le choix d'une bijection $I \simeq [1,n]$,
indifférent grâce à la stabilité par permutation ; les axiomes se transcrivent
en des énoncés (M$'$1) à (M$'$4) pour les applications $u : I \to J$ entre
ensembles finis. C'est sous cette forme qu'ils servent dans tout le reste du
dossier.

\subsection*{Les cas triviaux}

Par (M1), si l'un des $\mathcal{M}_n$ est vide, tous le sont, et les axiomes
sont alors vérifiés trivialement ; on exclut ce cas :
\[ (*) \qquad \mathcal{M}_1 \neq \emptyset . \]
De même, si l'un des $\mathcal{M}_n$ contient $\emptyset$, tous le contiennent.
Sinon, chaque $\mathcal{M}_n$, non vide, stable par intersection et ne
contenant pas $\emptyset$, est une base de filtre sur $R^n$ — situation
étrangère à ce qu'on cherche. On l'exclut aussi :
\[ (**) \qquad \emptyset \in \mathcal{M}_1 , \]
d'où $\emptyset \in \mathcal{M}_n$ pour tout $n$. Ce n'est pas une vraie
restriction : si (**) manque, les $\mathcal{M}_n \cup \{\emptyset\}$ vérifient
encore (M1) à (M4). Enfin si $\mathcal{M}_n = \{\emptyset\}$ pour un $n$, c'est
vrai pour tout $n$, et l'on exclut ce dernier cas trivial :
\[ (***) \qquad \text{il existe } A \in \mathcal{M}_1 \text{ non vide.} \]

Sous (**) et (***), on peut admettre aussi l'ensemble d'indices vide, en posant
$\mathcal{M}_\emptyset = \mathfrak{P}(R^\emptyset) = \{\emptyset, \{e\}\}$, où
$R^\emptyset = \{e\}$ est un point : les propriétés (M$'$1) à (M$'$4) restent
vraies, et ce choix est d'ailleurs forcé, puisque l'image d'une partie
modérée non vide par $R^I \to R^\emptyset$ est $\{e\}$.\footnote{La page écrit
« si on pose pour $I \neq \emptyset$ », lecture douteuse, puis
$\mathcal{M}_\emptyset = \ldots$ ; c'est le cas $I = \emptyset$ qui est
traité. La vérification de (M$'$1) à (M$'$4) est laissée sur la page à des
points de suspension ; elle est immédiate, et c'est (**) qui la rend vraie.}

\subsection*{Les points modérés}

Pour $x = (x_1, \dots, x_n) \in R^n$, on a $\{x\} \in \mathcal{M}_n$ si et
seulement si $\{x_i\} \in \mathcal{M}_1$ pour tout $i$ : dans un sens par
projection (M1), dans l'autre par produit (M4).\footnote{La page écrit
« $\{x\} \in R^n$ » à gauche de l'équivalence ; c'est $\{x\} \in
\mathcal{M}_n$ qui est entendu.} On pose donc
\[ R_0 = \{x \in R \mid \{x\} \in \mathcal{M}_1\} , \]
et l'ensemble des points modérés de $R^I$ est $R_0^I$. Si $R_0$ a au moins
deux points $x \neq y$, alors $\{x\} \cap \{y\} = \emptyset$ est modérée, et
(**) est automatique.

Dans le cas qui guide l'intuition — $R = \mathbb{R}$, parties semi-algébriques
compactes définies par des polynômes à coefficients rationnels — $R_0$ serait
le corps des nombres réels algébriques : les points modérés sont les points
« définissables » de la structure.\footnote{Exemple de l'édition ; le dossier
n'en donne aucun de cette sorte. Il annonce à la page 20 que $R_0$ est un
sous-anneau, et à la page 22 un sous-corps, ce qui va dans ce sens.}

\subsection*{Exemples (page 5)}

En dehors des deux cas triviaux, on a les exemples : a) les parties de
cardinal au plus $1$ ; b) les parties finies ; c) toutes les parties. Tous
vérifient (M1) à (M4), et dans les trois $R_0 = R$. « Aucun de ces
[exemples n'est intéressant] pour nous ! »\footnote{La phrase est coupée par un
long ajout interlinéaire en biais, dont on ne lit que des fragments
(« exemples », « discrets ») ; la fin de la phrase de la page est elle-même en
partie illisible, et la restitution entre crochets est de l'édition. Une note
marginale reliée à l'exemple a) renvoie aux axiomes et à « l'exclusion de
(2) » ; sa lecture est trop incertaine pour qu'on en tire quoi que ce soit.}
Ce sont les exemples où la projection ne crée rien : la structure ne voit pas
la géométrie de $R$.

\pagerange{5}{7}
\section*{(II) Espaces modérés (pages 5 à 7)}

\subsection*{Applications modérées}

Pour $A \in \mathcal{M}_I$ et $B \in \mathcal{M}_J$, une application $f : A \to
B$ est dite \emph{modérée} si son graphe $\Gamma_f \subset A \times B \subset
R^{I \sqcup J}$ est une partie modérée.\footnote{La page écrit l'indice
« $I \cup J$ » ; c'est la réunion disjointe qui est entendue, pour que
$R^I \times R^J \simeq R^{I \sqcup J}$.}

\textbf{Proposition.} a) $\mathrm{id}_A$ est modérée ; b) le composé de deux
applications modérées est modéré ; c) l'inverse d'une bijection modérée est
modéré.

En effet, le graphe de $\mathrm{id}_A$ est l'image de $A$ par un plongement
diagonal (M1). Celui de $g \circ f$, pour $f : A \to B$ et $g : B \to C$, est la
projection sur $R^{I \sqcup K}$ de $(\Gamma_f \times C) \cap (A \times
\Gamma_g)$, partie modérée de $R^{I \sqcup J \sqcup K}$ par (M4), (M3) et
(M1). Celui de $f^{-1}$ se déduit de $\Gamma_f$ par une permutation des
coordonnées.\footnote{En face de la Proposition, la marge porte « via M2 ».
La vérification ci-dessus n'emploie que (M1), (M3) et (M4) ; (M2) sert à la
réduction de (M3) et à l'intersection, ce que la marge vise peut-être.} Plus
généralement, l'image d'une partie modérée par une application modérée est
modérée, et l'image réciproque aussi : $f(A') = \mathrm{pr}_J\bigl(\Gamma_f
\cap (A' \times R^J)\bigr)$ et $f^{-1}(B') = \mathrm{pr}_I\bigl(\Gamma_f \cap
(R^I \times B')\bigr)$, par (M3) puis (M1).\footnote{Ces deux formules
explicitent une note marginale de la page 7, à demi lisible : « Introduire les
parties modérées des espaces modérés, \ldots\ les parties $f(A)$,
$f^{-1}(B)$ ».}

\subsection*{La catégorie des espaces modérés}

On obtient la catégorie $\mathrm{Mod}_{\mathcal{M}_*}$ des \emph{modèles
modérés} — les couples $(I, A)$ avec $A \in \mathcal{M}_I$ — et un foncteur
$(I, A) \mapsto A$ vers les ensembles, fidèle, et conservatif par c). On la
remplace par une catégorie équivalente $\Sigma = \mathrm{Esp}_{\mathcal{M}_*}$
d'\emph{espaces modérés} : un ensemble muni d'une structure transportée d'un
modèle par une bijection, deux telles bijections définissant la même structure
quand elles diffèrent d'une bijection modérée. Le foncteur d'oubli $X \mapsto
|X|$ est alors fidèle, conservatif et de plus \emph{transportable} : pour tout
objet $X$ et toute bijection $\alpha : |X| \xrightarrow{\sim} E$, il existe
une structure modérée sur $E$, et une seule, faisant de $\alpha$ un
isomorphisme de $\Sigma$ — c'est le transport de structure de
Bourbaki.\footnote{Le terme « transportable » est celui de la page.
Aujourd'hui on dirait que le foncteur d'oubli est \emph{uniquement
transportable} (Adámek, Herrlich et Strecker), ce qui fait de $\Sigma$, à
équivalence près, une catégorie concrète d'ensembles structurés. La page
écrit « $\exists$ un unique $f : X \to Y$ qui relève $\alpha$ » : c'est le
couple formé de $Y$ et de $f$ qui est unique, $|Y| = E$ étant fixé.}

La marge note qu'on n'exclut pas les modèles où $R_0$ est réduit à un point.

\subsection*{Produits et limites finies (page 6)}

\textbf{Proposition.} La catégorie $\Sigma$ a des produits de deux objets, et
le foncteur d'oubli les respecte. Elle a aussi des noyaux de couples, respectés
par l'oubli ; donc, dès qu'elle a un objet final, toutes les limites
projectives finies, que l'oubli respecte.

Le produit de $A \in \mathcal{M}_I$ et $B \in \mathcal{M}_J$ est $A \times B
\in \mathcal{M}_{I \sqcup J}$ (M4) ; ses projections sont modérées (M1), et
pour $f : C \to A$, $g : C \to B$ modérées, le graphe de $(f,g)$ est
$(\Gamma_f \times B) \cap \mathrm{pr}_{C,B}^{-1}(\Gamma_g)$, modéré par (M3).
Le noyau de $f, g : A \rightrightarrows B$ est la projection sur $R^I$ de
$\Gamma_f \cap \Gamma_g$, modérée par (M1) — éventuellement vide, d'où l'utilité
de (**).

Pour l'objet final, la page introduit
\[ (\mathrm{M}5) \qquad R_0 \neq \emptyset , \]
et un point modéré $\{x\}$, $x \in R_0$, est alors final : le graphe de
l'application constante $A \to \{x\}$ est $A \times \{x\}$.\footnote{Deux
remarques. D'abord, si l'on admet l'ensemble d'indices vide, comme la page 4
le permet et comme la marge de la page 5 le dit expressément (« on n'exclut
pas $I$ vide »), le point $R^\emptyset = \{e\}$ est déjà un objet final, sans
(M5). (M5) n'est nécessaire que si l'on se borne aux $I$ non vides. Ensuite,
sur la page, tout le passage qui va de « le foncteur $X \mapsto |X|$ » à
« limites finies » est encadré et barré de trois traits obliques, avec en
biais un ajout dont on lit « existence d'un objet final \ldots\ des produits
\ldots\ commute » : il a repris cet énoncé, et la page 8 le reformule sous
l'hypothèse $\operatorname{card} R_0 \geqslant 2$. On donne ici ce qui est vrai.}

\subsection*{La droite comme objet (pages 6 et 7)}

Si $R^n \in \mathcal{M}_n$ pour un $n \geqslant 1$, alors $R \in
\mathcal{M}_1$ par projection, et $R^I \in \mathcal{M}_I$ pour tout $I$ fini
par produit. D'où l'axiome, facultatif :
\[ (\mathrm{M}0) \qquad R \in \mathcal{M}_1 . \]
Sous (M0), $R$ est un objet de $\Sigma$, et tout objet de $\Sigma$ est
isomorphe à un sous-objet d'une puissance $R^n$ (produit dans $\Sigma$). Les
$\mathcal{M}_n$ se relisent alors dans $\Sigma$ : ce sont exactement les
sous-objets de $R^n$. En effet un monomorphisme de $\Sigma$ est injectif — son
couple noyau, qui existe et se calcule comme dans les ensembles, est trivial —
et une injection modérée $X \to R^n$ est un isomorphisme sur son image, qui est
modérée, par le c) de la Proposition de la page 5.

La page conclut que la donnée de $\mathcal{M}_*$ équivaut, à isomorphisme près,
à celle du triple $(\Sigma, R, X \mapsto |X|)$ soumis à des conditions :
a) $\Sigma$ a des limites projectives finies ; b) le foncteur d'oubli
\ldots ; c) tout objet de $\Sigma$ est un sous-objet d'une puissance de $R$.
Ce qui est établi, c'est un sens : $\mathcal{M}_*$ se reconstruit à partir du
triple, comme ensemble des images des monomorphismes $X \hookrightarrow
R^n$.\footnote{La condition b) est illégible sur la page (trois mots). On
attend « fidèle, conservatif, transportable, et commute aux limites
projectives finies », mais c'est une restitution. La réciproque — qu'un triple
vérifiant a), b), c) provienne d'un $\mathcal{M}_*$ — demande au moins que
l'image d'un sous-objet de $R^m$ par une application de coordonnées $R^m \to
R^n$ soit encore un sous-objet, c'est-à-dire une factorisation image dans
$\Sigma$ ; rien sur la page ne la fournit. Le dossier 121, à la page 58, se donne
explicitement de telles factorisations. Le titre final de la page, « Parlons aussi des opérations
\ldots\ modérées d'espaces modérés \ldots\ dans des $R^I$ », est trop lacunaire
pour être restitué.}

\pagerange{7}{13}
\section*{(III) Limites inductives finies (pages 7 à 13)}

\subsection*{Le problème}

La section annonce « le problème des quotients des espaces modérés, et des
limites inductives finies ».\footnote{L'introduction de la section, au bas
de la page 7, est presque entièrement illisible ; seuls ces mots sont sûrs.}
C'est la question de fond du dossier. Les limites projectives viennent
gratuitement des axiomes, parce que produits et projections y sont ; les
limites inductives — recoller, écraser — demandent de trouver dans un $R^n$ une
place où loger le résultat, et rien dans (M1)--(M4) ne la fournit.

\subsection*{Ensembles finis, et le principe de la page 8}

\textbf{Proposition (page 8).} Si $\operatorname{card} R_0 \geqslant 2$,
$\Sigma$ a des limites projectives finies, que l'oubli respecte. — Ce n'est
qu'une reprise de la page 6 : deux points modérés distincts donnent un objet
final et assurent (**).\footnote{Les quatre premières lignes de l'énoncé sont
biffées et surchargées ; seule la reprise finale est sûre dans son sens. Le
« moyennant (M5) » qui y figure est sans objet une fois $\operatorname{card}
R_0 \geqslant 2$ supposé.}

\textbf{Corollaire.} Sur tout ensemble fini, il existe une et une seule
structure modérée pour laquelle toutes les parties sont modérées : la
structure modérée discrète.\footnote{Ce corollaire demande la stabilité par
réunions finies, (M5$'$), qui n'apparaît qu'à la page suivante. Sans elle il
est faux : dans l'exemple a) de la page 5 (parties de cardinal au plus $1$),
$R_0 = R$ a autant de points qu'on veut, et pourtant aucun ensemble à deux
éléments ne porte de structure modérée. Avec (M5$'$) et $\operatorname{card}
R_0 \geqslant 2$, un ensemble à $n$ éléments s'envoie injectivement dans
$R_0^k$ dès que $2^k \geqslant n$, son image est une réunion finie de points
modérés, et toute application entre deux tels ensembles a un graphe fini de
points modérés, donc modéré : d'où existence et unicité.}

Le but annoncé est de construire les sommes amalgamées $X_1 \sqcup_Z X_2$ sans
supposer $Z$ fini. Le principe dont tout dépend est le suivant.

\textbf{Proposition (pages 8 et 11).} Soit $(X_i)_{i \in I}$ un diagramme fini
d'espaces modérés, et $X$ sa limite inductive dans les ensembles, avec les
applications canoniques $g_i : X_i \to X$. S'il existe sur $X$ une structure
modérée rendant les $g_i$ modérées, alors $X$, muni de cette structure, est la
limite inductive du diagramme dans $\Sigma$, et l'oubli la respecte.

\emph{Démonstration} (page 11). Un cône $(f_i : X_i \to Y)$ de $\Sigma$ définit
une application $f : X \to Y$ avec $f \circ g_i = f_i$ ; il faut voir qu'elle
est modérée. Comme $X$ est réunion des images des $X_i$, le graphe $\Gamma_f$
est la réunion des images des $\Gamma_{f_i}$ par les applications modérées $g_i
\times \mathrm{id}_Y$ ; ces images sont modérées, et leur réunion finie
aussi.\footnote{La dernière étape demande (M5$'$), stabilité par réunions
finies, que la page n'introduit qu'à la page 9, après l'énoncé de la page 8 et
avant la démonstration de la page 11 : l'énoncé est donc à lire sous (M5$'$).
Pour un conoyau de couple, $X$ est l'image d'un seul $X_i$ et la réunion est
superflue. Le début de la page 11 est d'une lecture incertaine ; la
démonstration est restituée d'après ce qui en est lisible, qui en donne
toutes les étapes.}

\subsection*{Réunions (page 9)}

On ajoute l'axiome :
\[ (\mathrm{M}5') \qquad \text{chaque } \mathcal{M}_n \text{ est stable par
réunions finies} \]
(y compris la réunion vide : $\emptyset \in \mathcal{M}_n$).\footnote{Sur la
page, l'étiquette est (M5), déjà employée page 6 pour $R_0 \neq \emptyset$ ;
voir la section des notations.} Alors : les parties modérées d'un espace
modéré sont stables par réunions finies ; $\emptyset$ porte une unique structure
modérée, qui en fait l'objet initial de $\Sigma$ ; et si $X$ est un espace
modéré réunion d'une famille finie de parties modérées $X_i$, $X$ est la somme
amalgamée des $X_i$ le long des $X_i \cap X_j$ : une application $f : X \to Y$
est modérée si et seulement si ses restrictions aux $X_i$ le sont, puisque
$\Gamma_f = \bigcup_i \Gamma_{f|X_i}$.

\textbf{Question.} Étant donnés deux plongements modérés $i_1 : Z
\hookrightarrow X_1$ et $i_2 : Z \hookrightarrow X_2$ dans $\Sigma$,

\begin{tikzcd}[column sep=small]
 & Z \arrow[dl, hook'] \arrow[dr, hook] & \\
X_1 & & X_2
\end{tikzcd}

la somme amalgamée $X_1 \sqcup_Z X_2$ existe-t-elle ? Pour $Z = \emptyset$,
c'est la question des sommes finies.\footnote{La page écrit « Si $Y =
\emptyset$ » : la lettre $Z$ du diagramme est écrite par-dessus une autre,
sans doute $Y$, restée dans la phrase. Pour les sommes finies, deux points
modérés distincts $a \neq b$ et (M5$'$) suffisent : on loge $X_1 \subset R^I$
et $X_2 \subset R^J$ dans $R^I \times R^J \times R$ comme $X_1 \times \{c_J\}
\times \{a\}$ et $\{c_I\} \times X_2 \times \{b\}$, avec $c_I, c_J$ des points
modérés.}

\subsection*{Le critère (pages 10 et 11)}

\textbf{Proposition.} Pour que la somme amalgamée de $X_1 \hookleftarrow Z
\hookrightarrow X_2$ existe dans $\Sigma$ et que l'oubli la respecte, il faut
et il suffit qu'il existe une famille finie $(f_\alpha, g_\alpha)_{\alpha \in
A}$ de couples de fonctions \emph{modérées} $f_\alpha : X_1 \to R$, $g_\alpha
: X_2 \to R$ telles que :

\begin{itemize}
\item[a)] $f_\alpha \circ i_1 = g_\alpha \circ i_2$ sur $Z$, pour tout $\alpha$ ;
\item[b)] les $f_\alpha$ séparent les points de $X_1$ et les $g_\alpha$ ceux de
  $X_2$, i.e.\ $f = (f_\alpha) : X_1 \to R^A$ et $g = (g_\alpha) : X_2 \to
  R^A$ sont injectives ;
\item[c)] pour $x_1 \in X_1 \smallsetminus i_1(Z)$ et $x_2 \in X_2 \smallsetminus
  i_2(Z)$, il existe $\alpha$ avec $f_\alpha(x_1) \neq g_\alpha(x_2)$.\footnote{Le mot « modérées » n'est pas dans l'énoncé de la page, mais il l'est
dans le NB qui le généralise, et le critère est faux sans lui. En b), la page
écrit « $X_1 \to R^I$ » pour le premier plongement ; l'ensemble d'indices est
$A$ dans tout le reste.}
\end{itemize}

Sous (M5$'$) — et c'est la seule chose à ajouter —, la preuve est la suivante.
Les conditions a), b), c) disent exactement que $f$ et $g$ se recollent en une
\emph{injection} $F$ de la somme amalgamée ensembliste $X = X_1 \sqcup_Z X_2$
dans $R^A$ : b) sépare les points d'un même côté, c) sépare un point de $X_1$
hors de $Z$ d'un point de $X_2$ hors de $Z$, et un point $i_1(z)$ d'un point
$x_2 \neq i_2(z)$ est séparé par $g$ grâce à a) et b). L'image $F(X) = f(X_1)
\cup g(X_2)$ est modérée, on transporte sa structure à $X$, les $g_i$ sont
modérées puisque $F \circ g_1 = f$ et $F \circ g_2 = g$ le sont, et la
Proposition des pages 8 et 11 conclut. Inversement, si la somme existe et que
l'oubli la respecte, un plongement $F : X \hookrightarrow R^A$ donne $f = F
\circ g_1$ et $g = F \circ g_2$.

\textbf{NB.} De même, pour un diagramme fini quelconque, la structure modérée
quotient existe sur la limite inductive ensembliste $X$ si et seulement s'il
existe une famille finie de systèmes compatibles $(f_{\alpha i} : X_i \to
R)_{i}$ de fonctions modérées, définissant des $f_\alpha : X \to R$ qui
séparent les points de $X$.

\textbf{Corollaire (page 11).} Soit $f : X \to Y$ une application modérée et
$\mathcal{R}_f = X \times_Y X$ la relation d'équivalence qu'elle définit —
partie modérée de $X \times X$, et relation d'équivalence dans $\Sigma$. Le
quotient $X / \mathcal{R}_f$ existe dans $\Sigma$ et s'identifie
canoniquement à $f(X)$ muni de sa structure induite. Autrement dit, dans
$\Sigma$, toute application modérée se factorise en un épimorphisme effectif
suivi d'un plongement, comme dans les ensembles.\footnote{La traduction en
termes de factorisation est de l'édition. Le corollaire suit de la Proposition
des pages 8 et 11 appliquée au conoyau de $\mathcal{R}_f \rightrightarrows X$,
sans même (M5$'$). Le « Cor.\ 2 » qui suit sur la page — de nouveau le
diagramme $X_1 \hookleftarrow Z \hookrightarrow X_2$ — est barré de quatre
traits ; il est repris page 12.}

\subsection*{Deux axiomes (page 12)}

Le critère dit ce qu'il faut ; deux axiomes le rendent applicable.

\[ (\mathrm{M}6) \quad
\begin{cases}
\text{pour } Z \subset X \text{ dans } \mathcal{M}_n, \text{ il existe une famille
finie } (f_\alpha)_{\alpha \in A} \\
\text{de fonctions modérées } X \to R \text{ et des } e_\alpha \in R_0
\text{ tels que } Z = \bigcap_\alpha f_\alpha^{-1}(e_\alpha) ; \\
\text{de plus, on peut supposer que les } f_\alpha \text{ séparent les points de }
X \smallsetminus Z .
\end{cases} \]

Avec $F = (f_\alpha) : X \to R^A$ et $e = (e_\alpha)$, $Z = F^{-1}(e)$, et $F$
est injective hors de $Z$ : $F$ identifie $X/Z$ — $X$ où $Z$ est écrasé en un
point — à la partie modérée $F(X)$ de $R^A$. L'axiome (M6) signifie donc que
le quotient $X/Z$ existe dans $\Sigma$ « au sens strict », c'est-à-dire que
l'oubli le respecte.\footnote{Pour $Z$ non vide ; pour $Z = \emptyset$, $X/Z$
est $X$ augmenté d'un point. Que $e_\alpha$ soit pris dans $R_0$ est
essentiel : c'est ce qui rend $\{e\}$ modéré, et l'on s'en sert page 13. La
seconde moitié de l'axiome (« de plus, on peut supposer », écrit « OPS ») n'est
pas une conséquence de la première dans le cadre ensembliste ; elle le devient
dès que $R$ est un corps, et c'est ce que prouvent les pages 22 et 23. Sur la
page, sous (M6), un diagramme carré $Z \hookrightarrow R^n \times \{e\}$,
$X \hookrightarrow R^n \times R^A$ est pris dans un bloc barré ; il est repris
correctement page 13.}

\[ (\mathrm{M}7) \quad
\begin{cases}
\text{pour } Z \subset X \text{ dans } \mathcal{M}_n \text{ et } f : Z \to R
\text{ modérée,} \\
\text{il existe } g : X \to R \text{ modérée qui prolonge } f .
\end{cases} \]

C'est un théorème de Tietze modéré.\footnote{Dans le cadre o-minimal, les deux
axiomes sont des théorèmes : un fermé définissable est le lieu des zéros
d'une fonction définissable continue, et les fonctions définissables continues
sur un fermé se prolongent (théorème de Tietze définissable ; voir le livre de
van den Dries, \emph{Tame topology and o-minimal structures}, 1998). Cela
suppose la topologie et la structure de corps ordonné, qu'ici le dossier n'a
pas encore introduites ; ce sont des rapprochements de l'édition.}

\subsection*{Existence des sommes amalgamées (pages 12 et 13)}

\textbf{Proposition.} Sous (M5$'$), (M6) et (M7), les sommes amalgamées $X_1
\sqcup_Z X_2$ le long de plongements $Z \hookrightarrow X_1$, $Z
\hookrightarrow X_2$ existent dans $\Sigma$, et l'oubli les respecte.

\emph{Démonstration} (page 13). Soit $\varphi : Z \hookrightarrow R^I$ un
plongement modéré. Par (M7), coordonnée par coordonnée, il se prolonge en
$\varphi_i : X_i \to R^I$ ($i = 1, 2$). Par (M6) appliqué à $Z_i = u_i(Z)
\subset X_i$, on a $F_i : X_i \to R^{A_i}$ avec $Z_i = F_i^{-1}(e_i)$, $e_i \in
R_0^{A_i}$, injective hors de $Z_i$. Alors $(\varphi_1, F_1) : X_1 \to R^I
\times R^{A_1}$ est un plongement — les points hors de $Z_1$ sont séparés par
$F_1$, ceux de $Z_1$ par $\varphi$ — et le carré

\begin{tikzcd}[column sep=small, row sep=small]
X_1 \arrow[r, hook] & R^I \times R^{A_1} \\
Z \arrow[u, hook] \arrow[r, hook] & R^I \times \lbrace e_1 \rbrace
\arrow[u, hook]
\end{tikzcd}

est cartésien.\footnote{Dans ce diagramme, la page trace la flèche verticale
de gauche vers le bas, de $X_1$ vers $Z$ ; on la donne dans le sens de
l'inclusion, qui est celui de la flèche de droite.} De même pour $X_2$. On
range alors les deux morceaux à des étages différents d'un même espace :

\begin{tikzcd}[column sep=small, row sep=small, nodes={font=\scriptsize}]
X_1 \arrow[r, hook] & R^I \times R^{A_1} \times \lbrace e_2 \rbrace
\arrow[r, hook] & R^I \times R^{A_1} \times R^{A_2} \\
 & R^I \times \lbrace e_1 \rbrace \times \lbrace e_2 \rbrace \arrow[u, hook]
\arrow[r, hook] & R^I \times \lbrace e_1 \rbrace \times R^{A_2} \arrow[u,
hook] \\
Z \arrow[uu, hook] \arrow[rr, hook] & & X_2 \arrow[u, hook]
\end{tikzcd}

C'est le diagramme marqué (*) en marge de la page 13.\footnote{Il le
reprend d'une première version barrée. La page trace la flèche de la première
ligne de droite à gauche, sans crochet, et note les crochets de la colonne de
droite par de petites marques de sens incertain ; on les donne toutes dans le
sens des inclusions, qui est le seul où le diagramme commute.} $X_1$ est logé
dans l'étage $R^{A_2}$-coordonnée $= e_2$, $X_2$ dans l'étage
$R^{A_1}$-coordonnée $= e_1$. Un point commun aux deux images a ses deux
coordonnées égales à $e_1$ et $e_2$, donc provient de $Z_1$ et de $Z_2$, et la
coordonnée $R^I$ — qui est $\varphi$ sur $Z$ — dit qu'il provient du même point
de $Z$. Les deux images se rencontrent donc exactement le long de l'image de
$Z$, et leur réunion, modérée par (M5$'$), est en bijection avec la somme
amalgamée ensembliste ; c'est le critère des pages 10 et 11, et la Proposition
des pages 8 et 11 conclut.\footnote{La page s'arrête après le diagramme ; la
conclusion est restituée. Les étages $\{e_1\}$, $\{e_2\}$ sont des parties
modérées parce que $e_1, e_2$ sont des points modérés, ce que (M6) exige.}

\pagerange{14}{16}
\section*{(IV) Ind-structures modérées strictes (pages 14 à 16)}

Un \emph{ind-espace modéré strict} est un ensemble $E$ muni d'un ensemble
$\mathcal{M}_E$ de parties — ses parties modérées —, stable par réunions finies
et recouvrant $E$, chacune munie d'une structure modérée, ces structures
s'induisant mutuellement, et tel que $A \in \mathcal{M}_E$, $B$ partie modérée
de $A$ entraînent $B \in \mathcal{M}_E$. Il revient au même de se donner la
notion d'application modérée d'un espace modéré dans $E$. C'est un
\emph{ind-objet strict} de $\Sigma$ : une réunion filtrante d'espaces modérés
le long de plongements.\footnote{Le titre porte en interligne « (ou
Ind-espaces modérés) », avec une majuscule nette ; les « ind- » minuscules
des pages suivantes sont le même préfixe. L'expression « ind-objet strict »
est de l'édition ; c'est la notion de SGA~4 restreinte aux systèmes de
monomorphismes.}

On suppose
\[ (\mathrm{M}8) \qquad R \text{ est réunion de ses parties modérées} ; \]
$R$, et les $R^I$ pour $I$ fini, sont alors canoniquement des ind-espaces
modérés stricts.\footnote{La page ajoute au-dessus de la ligne « (impliqué par
M5) », lecture douteuse. Ni l'une ni l'autre de ses deux (M5) n'entraîne (M8) :
si $R_0 \subsetneq R$ et que $\mathcal{M}_n$ est l'ensemble des parties finies
de $R_0^n$, les axiomes (M1) à (M4), (M5) et (M5$'$) sont vérifiés et $R$ n'est
pas réunion de ses parties modérées. (M8) est en revanche conséquence de
(M0).}

\textbf{Proposition.} Sous (M1) à (M5$'$), la catégorie $\Sigma^{\mathrm{ind}}$
des ind-espaces modérés stricts a des limites projectives finies, que l'oubli
respecte ; elle a aussi des sommes finies, respectées par l'oubli.\footnote{La
marge porte « M1 : M5 ». Le sigle de la catégorie est $\mathrm{Esp}$ avec un
indice surchargé, $\mathcal{M}_*$ ou $\mathcal{M}$ ; on écrit
$\Sigma^{\mathrm{ind}}$ pour la distinguer de $\Sigma$. La page met les sommes
finies sous l'hypothèse « $R_0 \neq \emptyset$ ». La construction qu'on connaît
(note de la page 9) demande deux points modérés distincts, pour séparer les
deux termes ; avec un seul, l'édition ne voit pas de preuve. Une remarque
marginale, en partie lisible, ajoute « stricts \ldots\ (en principe !) ».}

\subsection*{Parties localement modérées (pages 15 et 16)}

Une partie $\mathfrak{Y}$ d'un ind-espace $\mathfrak{X}$ est \emph{localement
modérée} si $A \cap \mathfrak{Y}$ est modérée pour toute partie modérée $A$ de
$\mathfrak{X}$. Ces parties sont stables par intersections finies — la
famille vide donnant $\mathfrak{X}$ lui-même — et par image réciproque par les
morphismes d'ind-espaces. Une application modérée d'un espace modéré dans un
$R^I$ est un cas particulier de morphisme d'ind-espaces.

Sous (M6) et (M7), la page annonce l'existence, dans $\Sigma^{\mathrm{ind}}$ et
au sens strict, des quotients $\mathfrak{X}/\mathfrak{Y}$ pour $\mathfrak{Y}$
localement modérée, et des sommes amalgamées $\mathfrak{X}_1
\sqcup_{\mathfrak{Z}} \mathfrak{X}_2$ le long d'immersions correspondant à des
parties localement modérées $\mathfrak{Z}_i \subset \mathfrak{X}_i$, avec
$\mathfrak{Z} \simeq \mathfrak{Z}_1 \simeq \mathfrak{Z}_2$. Aucune démonstration
n'est esquissée : on passerait à la limite sur les parties modérées dans la
construction des pages 12 et 13.\footnote{La page écrit « dans
$\mathrm{Esp}_{\mathcal{M}_*}$ » ; c'est la catégorie des ind-espaces qui est
entendue, puisque $\mathfrak{X}$ en est un objet. La dernière phrase, sur la
méthode, est de l'édition.}

\textbf{NB.} Contrairement à ce qui se passe dans $\Sigma$, où les sous-objets
d'un objet sont exactement ses parties modérées (page 7), un sous-objet d'un
ind-espace ne correspond pas nécessairement à une partie localement
modérée.\footnote{Exemple de l'édition, dans le cadre semi-algébrique
compact : l'ensemble $\mathbb{N}$, ind-espace discret réunion de ses parties
finies, s'envoie injectivement dans la droite par $n \mapsto 1/(n+1)$ ; l'image
coupe $[0,1]$ en un ensemble infini dénombrable, qui n'est pas semi-algébrique.
La page porte peut-être sur les deux sigles $\mathrm{Esp}$ une marque qui
distinguerait la catégorie des ind-espaces.}

\pagerange{16}{19}
\section*{(V) Relations avec les topologies (pages 16 à 19)}

\subsection*{Parties modérées compactes}

Supposons $R$ muni d'une topologie ; les $R^I$ et les modèles modérés en
héritent. On veut que les applications modérées soient continues, et il suffit
pour cela que les parties modérées soient compactes, d'où l'axiome
\[ (\mathrm{M}9) \qquad R \text{ est séparé, et les parties modérées des }
R^n \text{ sont compactes.} \]
En effet, une application $f : A \to B$ entre compacts séparés dont le graphe
est fermé — ici compact — est continue. Le foncteur d'oubli se factorise alors
par la catégorie des espaces compacts séparés,
\[ \Sigma \longrightarrow (\text{espaces compacts}) , \]
et ce foncteur est encore fidèle, conservatif — une bijection continue entre
compacts séparés est un homéomorphisme — et transportable. Réciproquement,
quand $R \in \mathcal{M}_1$ (donc $R$ compact), la page annonce une
caractérisation des $\mathcal{M}_*$ par un tel foncteur et un objet $R$ dans
une puissance duquel tout objet se plonge, analogue à celle de la page
7.\footnote{Sous (M0) et (M9), $R$ est compact ; le dossier ne revient pas
sur la tension avec la section (VI), où $R$ est un corps : un corps
topologique séparé et compact est fini. Dans la section (VI) il a d'ailleurs
« oublié les structures topologiques ».}

\subsection*{Ind-espaces et topologie limite inductive (page 17)}

Tout ind-espace modéré est muni de la topologie limite inductive de ses
parties modérées. Celle-ci induit sur chaque partie modérée sa topologie
propre, et chaque partie modérée y est fermée.\footnote{La page met une
hypothèse : l'espace serait « ind-dénombrable », limite inductive de
\emph{suites} croissantes d'espaces modérés (« dénombrable » est une lecture
douteuse). Pour l'énoncé tel qu'il est écrit, elle est superflue : si $F$ est
fermé dans une partie modérée $A$, alors pour toute partie modérée $B$, $F$
est compact dans le compact séparé $A \cup B$, donc fermé, et $F \cap B$ est
fermé dans $B$ ; $F$ est donc fermé dans l'ind-espace, ce qui donne les deux
assertions. Là où la dénombrabilité sert vraiment, c'est pour que la topologie
limite inductive commute aux produits finis : c'est le cas pour les réunions
de suites croissantes de compacts (les « $k_\omega$-espaces »), et non en
général. L'édition ne sait pas laquelle des deux choses la page avait en
vue.}

\subsection*{Voisinages modérés (pages 17 et 18)}

On veut que dans un espace modéré tout point ait un système fondamental de
voisinages modérés ; il suffit que ce soit vrai dans $R$, car alors c'est vrai
dans les $R^n$ (par produits) et dans leurs parties modérées (par
intersection). D'où l'axiome
\[ (\mathrm{M}10) \qquad \text{dans } R, \text{ tout point a un système
fondamental de voisinages appartenant à } \mathcal{M}_1 . \]
Sous (M9), (M10) entraîne que $R$ est localement compact.

On dispose alors de la notion d'application \emph{modérée au voisinage d'un
point}.

\textbf{Proposition (page 18).} Sous (M5$'$), (M9) et (M10), une application
d'un espace modéré $X$ dans un ind-espace modéré est modérée si et seulement si
elle l'est au voisinage de chaque point.

En effet, $X$ étant compact, il est recouvert par un nombre fini de voisinages
modérés sur lesquels l'application est modérée, et son graphe est la réunion
finie des graphes des restrictions.\footnote{La page ne donne pas la preuve ni
la liste des axiomes ; on la donne, avec les hypothèses dont elle se sert.}
Les applications modérées de $X$ dans $Y$ sont donc les sections d'un
faisceau sur $X$, le \emph{faisceau des applications modérées}. La structure
modérée de $X$ est connue dès qu'on connaît les fonctions modérées $X \to R$
— c'est vrai déjà dans le cadre ensembliste, sous (M1)$_0$, (M2), (M3), (M4),
puisqu'une application dans $R^n$ est modérée si et seulement si ses
coordonnées le sont —, donc aussi dès qu'on connaît le faisceau des germes de
fonctions modérées.\footnote{« (M1)$_0$ » : l'indice $0$ est net sur la page,
et revient dans une note en biais de la page 7 que l'on ne sait pas lire. Le
dossier ne définit nulle part cette variante de (M1) ; on garde le sigle sans
lui attribuer de contenu.} Un espace modéré devient ainsi un espace annelé —
c'est le point de vue des « espaces définissables » de la théorie
o-minimale.\footnote{Rapprochement de l'édition.}

\subsection*{Espaces localement modérés (page 19)}

Un \emph{espace localement modéré} est un ind-espace modéré dont la topologie
est séparée, dont tout point a un voisinage modéré — de sorte qu'il est
localement compact — et même un système fondamental de voisinages modérés, et
dont toute partie compacte est contenue dans une partie modérée : les parties
modérées forment alors une partie cofinale, filtrante, de l'ensemble des
compacts.\footnote{La page dit « les parties compactes sont les parties
contenues dans une partie modérée » : lu à la lettre, toute partie d'une
partie modérée serait compacte, ce qui est faux. On lit l'inclusion dans le
seul sens vrai, que la suite de la phrase (« un ensemble filtrant cofinal :
celui des parties compactes ») confirme. Le mot « ind-modéré » est écrit
au-dessus de « loc.\ cpct », biffé.}

Par (M10), $R$ est localement modéré — un compact de $R$ est recouvert par un
nombre fini de voisinages modérés, dont la réunion est modérée par (M5$'$) — et
donc aussi les $R^I$. Les espaces localement modérés forment une
sous-catégorie pleine de $\Sigma^{\mathrm{ind}}$, stable par limites
projectives finies, ce que la page affirme sans le montrer.\footnote{Il note
cette sous-catégorie « $\mathrm{Esploc}_{\mathcal{M}_*}$ ». Une note en biais
au bas de la page est presque illisible.}

\pagerange{20}{23}
\section*{(VI) Relation avec une structure de corps (pages 20 à 23)}

\subsection*{L'axiome (M11)}

On oublie la topologie et l'on suppose que $R$ est un \emph{corps}. Outre les
axiomes qui font de $\Sigma$ une catégorie à limites projectives finies, on
demande :\footnote{La page dit « en plus de (M1) et (M4), minimum pour
construire $\mathrm{Esp}_{\mathcal{M}_*}$ avec ses $\varprojlim$ finies ».
Les limites finies de la page 6 se servent aussi de la stabilité par
intersection, donc de (M3), ou de (M2) avec (M4) ; un indice surchargé après
« (M1) » visait peut-être cela. Une première version de (M11), ne portant que
sur $a \mapsto -a$, est barrée de sept traits.}

\[ (\mathrm{M}11) \quad
\begin{cases}
\text{a) } 0 \neq 1 \text{ sont des points modérés : } \{0\}, \{1\} \in
\mathcal{M}_1 ; \\
\text{b) pour } A \in \mathcal{M}_2, \text{ les fonctions } (x,y) \mapsto x - y
\text{ et } (x,y) \mapsto xy \\
\qquad \text{sur } A \text{ sont modérées,} \\
\qquad \text{i.e.\ } \{(x, y, x-y)\} \text{ et } \{(x, y, xy)\},\ (x,y) \in A,
\text{ sont dans } \mathcal{M}_3 ; \\
\text{c) pour } A \in \mathcal{M}_1,\ A \subset R^*, \text{ il existe } B \in
\mathcal{M}_1 \text{ avec } A^{-1} \subset B .
\end{cases} \]

La marge dit deux fois, la première biffée, qu'on pourrait supposer $B =
A^{-1}$ ; la page 22 montre que c'est en effet une conséquence. Une note
marginale ajoute que, sous (M0), b) se réduit au cas $A = R^2$.\footnote{Lecture
partielle : « NB si \ldots\ $A \in \mathcal{M}_2$ \ldots\ (M11) pour $A =
R^2$ ». La réduction est juste : si $R^2$ est modérée et que les graphes de
$x - y$ et $xy$ sur $R^2$ le sont, leurs restrictions à $A$ le sont par (M3).}

\textbf{Proposition.} Pour un espace modéré $X$, les fonctions modérées $X \to
R$ forment un sous-anneau de l'anneau de toutes les fonctions sur $X$. En
particulier $R_0$ est un sous-anneau de $R$.

En effet, pour $f, g$ modérées, $(f, g) : X \to R^2$ est modérée, d'image $A
\in \mathcal{M}_2$, et $f - g$, $fg$ sont les composés de $(f,g)$ avec les
fonctions modérées de b) ; les constantes $0$ et $1$ sont modérées par a). Une
fonction constante de valeur $c$ est modérée si et seulement si $c \in R_0$ —
son graphe est $X \times \{c\}$ —, et $R_0$ est le sous-anneau des
constantes.\footnote{La page s'arrête au bas du feuillet sur « En particulier
$R_0$ est un sous-anneau de $R$, et les \ldots » ; la page 21 ne la continue
pas visiblement. La marge porte, verticalement, un fragment où l'on lit
« $R_0$ intègre », lecture douteuse — ce qui est vrai de tout sous-anneau d'un
corps.}

\subsection*{Inverses (pages 21 et 22)}

La page 21 s'ouvre sur un signe isolé, un D à double trait entouré d'une
boucle, suivi de « $=$ », dont le dossier ne dit pas la fonction.\footnote{On
ne lui donne pas de sens. La transcription le rend par $\mathbb{D}$.} Suit une
remarque : sous (M8), sur les ind-espaces $R^I$, les fonctions ind-modérées
contiennent en tout cas l'anneau $R_0[(X_i)_{i \in I}]$ des polynômes à
coefficients dans $R_0$ en les fonctions coordonnées — ce que donne la
Proposition précédente sur chaque partie modérée —, et les applications
$R_0$-polynomiales entre les $R^I$ sont modérées et transforment les parties
modérées en parties modérées.\footnote{Cette dernière clause est un ajout
interlinéaire serré, lu « Dans les applications $R_0$-polynomiales $R^I \to X$
modérées, et transforment des parties modérées en itou » ; le « $X$ » et
plusieurs mots sont douteux, et l'on restitue « entre les $R^I$ ». La page ne
dit pas si $I$ peut être infini ; partout ailleurs dans le dossier il est
fini, et on le garde fini. « Ensembles », souligné ou biffé, est suivi de
« anneaux » : on lit « anneaux ».}

\textbf{Corollaire.} Une fonction modérée $f$ sur un espace modéré $X$ est
inversible dans l'anneau des fonctions modérées si et seulement si elle ne
s'annule pas sur $X$.

Comme $f(X)$ est une partie modérée de $R^*$, cela revient au

\textbf{Corollaire.} Pour $A \in \mathcal{M}_1$, $A \subset R^*$, la restriction
de $x \mapsto x^{-1}$ à $A$ est modérée : $\{(x, x^{-1}) \mid x \in A\} \in
\mathcal{M}_2$.

En effet, cet ensemble est $\{(x, y) \in A \times B \mid xy = 1\}$, où $B \in
\mathcal{M}_1$ contient $A^{-1}$ (M11 c) : c'est l'image réciproque du point
modéré $1$ par la fonction modérée $(x,y) \mapsto xy$ sur $A \times B$. En
projetant sur la seconde coordonnée, $A^{-1}$ est elle-même modérée, ce que
suggérait la marge.

\textbf{Corollaire (page 22).} $R_0$ est un sous-corps de $R$ : pour $x \in
R_0^*$, appliquer ce qui précède à $A = \{x\}$.

\subsection*{Retour sur (M6) (pages 22 et 23)}

Soit $Y$ une partie modérée d'un espace modéré $X$, définie comme dans (M6)
par $Y = \bigcap_{\alpha} f_\alpha^{-1}(e_\alpha)$ avec $e_\alpha \in R_0$.
Quitte à remplacer $f_\alpha$ par $f_\alpha - e_\alpha$, qui est modérée parce
que la constante $e_\alpha$ l'est, on peut supposer $e_\alpha = 0$ : $Y$ est
le lieu des \emph{zéros communs} des $f_\alpha$.\footnote{La page écrit
« $e_\alpha \in R$ ». C'est $e_\alpha \in R_0$ qu'exige (M6), et c'est ce
qui rend la translation légitime : la constante $e_\alpha$ n'est modérée que si
$e_\alpha$ est un point modéré.} Soit $(g_\beta)_{\beta \in J}$ une famille
finie de fonctions modérées réalisant un plongement $X \hookrightarrow R^J$.
Les fonctions $f_\alpha$ et $f_\alpha g_\beta$ sont modérées, s'annulent sur
$Y$, ont $Y$ pour lieu des zéros communs, et \emph{séparent les points de $X
\smallsetminus Y$}. Soient en effet $x \neq y$ dans $X \smallsetminus Y$. Si
un $f_\alpha$ les sépare, c'est fini. Sinon $f_\alpha(x) = f_\alpha(y)$ pour
tout $\alpha$ ; comme $x \notin Y$, il existe $\alpha$ avec $\lambda =
f_\alpha(x) \neq 0$, et il existe $\beta$ avec $g_\beta(x) \neq g_\beta(y)$ ;
alors
\[ f_\alpha g_\beta(x) = \lambda\, g_\beta(x) \neq \lambda\, g_\beta(y) =
f_\alpha g_\beta(y) , \]
puisque $\lambda \neq 0$ dans le corps $R$.

\textbf{Proposition (page 23).} Sous (M1) à (M7) et (M11), toute partie modérée
$Y$ d'un espace modéré $X$ qui est le lieu des zéros communs d'une famille
finie de fonctions modérées est aussi celui d'une famille finie de fonctions
modérées séparant les points de $X \smallsetminus Y$ ; autrement dit $X/Y$ est
un espace modéré, au sens strict de la page 12.\footnote{La page écrit
« qui sépare les points de $X$ » : une telle famille, constante sur $Y$, ne
peut séparer les points de $Y$ dès que $Y$ en a deux ; c'est $X \smallsetminus
Y$ qui est entendu, comme dans l'argument qui précède. Les mots entre
« famille » et « $X/Y$ » sont illisibles, la parenthèse ouvrante manque ; la
lecture « $X/Y$ est un espace modéré » est sûre. Le « Dém : » qui précède
l'énoncé est resté vide : la démonstration est l'argument écrit juste
au-dessus, aux pages 22 et 23. « M1 : M7 » est peut-être « M1 à M7 ». Un mot
barré devant « modérée », lu « fermée », est douteux.}

Ainsi, dès que $R$ est un corps, la seconde moitié de (M6) est conséquence de
la première : c'est le seul endroit où le dossier allège ses propres axiomes.

\pagerange{24}{24}
\section*{À prouver (page 24)}

La dernière page n'est plus une construction mais un programme.

1) Pour un morphisme $X \to Y$ d'espaces modérés, il existerait une partie
modérée $Y'$ \emph{rare} dans $Y$ telle que, en notant $(U_i)_{i \in I}$ les
composantes connexes de $Y \smallsetminus Y'$ — leurs adhérences
$\overline{U}_i$ étant des parties modérées de $Y$ —, il existe pour chaque
$i$ un espace modéré $Z_i$ et une correspondance modérée
\[ R_i \subset (\overline{U}_i \times Z_i) \times X \]
induisant une bijection entre $U_i \times Z_i$ et la partie $X|U_i$ de $X$
au-dessus de $U_i$ : au-dessus de chaque $U_i$, $X$ est un fibré trivial au
sens modéré relatif.

Puis la question : un énoncé analogue vaut-il pour un \emph{diagramme fini}
d'espaces modérés au-dessus de $Y$ — est-il « trivial par morceaux »
au-dessus de $Y \smallsetminus Y'$ pour un $Y'$ modéré rare convenable ? « Ce
serait déjà intéressant » dans le cas d'un diagramme $Z \subset X \subset Y
\times [0,1]$ au-dessus de $Y$.\footnote{Le dernier mot avant « au dessus de
$Y$ » est illisible. Le crochet de $[0,1]$ est tracé par-dessus un premier
signe, peut-être un $I$.}

Ce programme suppose tacitement ce que le dossier n'a pas posé ensemble : une
topologie (« rare », « composantes connexes »), la finitude de $\pi_0(Y
\smallsetminus Y')$, et un ordre sur $R$ pour parler de $[0,1]$. Le choix
d'une \emph{correspondance} sur $\overline{U}_i$ plutôt que d'une bijection est
juste : une trivialisation sur $U_i$ ne se prolonge pas en bijection sur le
bord, mais son graphe a une adhérence.

C'est, à la lettre, le \emph{théorème de trivialité de Hardt} : établi pour les
applications semi-algébriques par Hardt en 1980, il vaut dans toute structure
o-minimale, et sous une forme compatible avec un nombre fini de parties — ce
qui répond par l'affirmative à la question du diagramme $Z \subset X$, la
partie rare $Y'$ étant de dimension strictement plus petite que
$Y$.\footnote{Voir le chapitre 9 du livre de van den Dries cité plus haut.
Ces pages, datées par l'inventaire « à partir de 1973 — à partir de 1978 »,
ne pouvaient connaître ni l'article de Hardt (1980) ni les structures
o-minimales (van den Dries, Pillay et Steinhorn, à partir de 1984–1986) ;
elles pouvaient connaître le théorème de Tarski–Seidenberg, les ensembles
semi-analytiques de Łojasiewicz (1964) et sous-analytiques de Hironaka
(1973), qu'elles ne citent pas. Rapprochement de l'édition, sans
affirmation de priorité.}

\end{document}
